
\paragraph{Input representation.}
Let $\mathbf{x}\in\mathbb{R}^{J\times F\times T}$ denote a root-centred motion sequence with $J{=}17$ Human3.6M joints~\citep{ionescu2014human36m}, $F{=}3$ Cartesian coordinates, and $T$ frames.  Let $f_y(\mathbf{x})\in[0,1]$ be the predicted probability assigned by a fixed classifier to class $y$.  For a set of interpretable players $N=\{1,\ldots,n\}$ and a coalition $S\subseteq N$, write $\mathbf{x}_S$ for the observed player components and $\bar S=N\setminus S$ for the hidden components.  MotionSHAP estimates the on-manifold value function
\begin{equation}
  v^{(\mathrm{on})}_{\mathbf{x}}(S)
  =
  \mathbb{E}_{p(\mathbf{x}'_{\bar S}\mid \mathbf{x}_S)}
  \left[
    f_y\!\left(\mathbf{x}_S\sqcup \mathbf{x}'_{\bar S}\right)
  \right],
  \label{eq:motion_von}
\end{equation}
where $\sqcup$ denotes splicing observed components from the target sequence with imputed hidden components.  We then use the same Kernel SHAP weighted least-squares estimator described in Sec.~3.2; methods differ only in the conditional distribution used to approximate Eq.~\eqref{eq:motion_von}.

\paragraph{Player sets.}
We evaluate three player granularities.  For temporal explanations, $n=K$ players are contiguous temporal windows, with $K\in\{4,8\}$.  In synthetic benchmarks these windows partition the sequence uniformly.  In real gait clips, the $K{=}4$ setting uses gait-phase windows: the walking direction is the first principal component of the pelvis trajectory, the stride period $P$ is the first peak of the normalised autocorrelation of $\max(\mathrm{RAnkle}_{\parallel},\mathrm{LAnkle}_{\parallel})-\mathrm{Pelvis}_{\parallel}$, and frames are assigned by phase $\varphi_t=(t\bmod P)/P$:
\begin{equation}
  W_k=\{t:\varphi_t\in[b_k,b_{k+1})\},
  \quad
  (b_0,b_1,b_2,b_3,b_4)=(0,0.30,0.60,0.80,1.00).
  \label{eq:windows}
\end{equation}
These windows correspond to initial contact/loading, mid-stance/terminal stance, pre-swing/initial swing, and mid-swing/terminal swing; when stride detection is unreliable we use uniform temporal quarters.  For spatial explanations, $n=J=17$ players correspond to individual joints.  Reported anatomical-group scores are obtained by summing joint-level Shapley values, preserving additivity.

\paragraph{Imputer interface.}
Both learned models below provide conditional completions for the same cooperative game.  Given a coalition $S$, we convert it to a binary observation mask $\mathbf{m}_S\in\{0,1\}^{J\times F\times T}$, where $m_{jft}=1$ means coordinate $(j,f,t)$ is observed.  The coalition value is estimated by Monte Carlo,
\begin{equation}
  \widehat{v}_{\mathbf{x}}(S)
  =
  \frac{1}{M}\sum_{r=1}^{M}
  f_y\!\left(
    \mathbf{x}_S\sqcup\hat{\mathbf{x}}_{\bar S}^{(r)}
  \right),
  \label{eq:vfn_mc}
\end{equation}
where each completion $\hat{\mathbf{x}}_{\bar S}^{(r)}$ is drawn from either MotionSHAP-VAEAC or MotionSHAP-Flow and observed coordinates are always overwritten with their original values before evaluating $f_y$.

% ------------------------------------------------------------------
\subsection{MotionSHAP-VAEAC}
\label{sec:architecture}
% ------------------------------------------------------------------

Our first imputer is a variational autoencoder with arbitrary conditioning~\citep{ivanov2018variational,olsen2022using}.  It represents a motion clip as a sequence of $T$ frame tokens, each containing the flattened pose in $\mathbb{R}^{JF}$ and the corresponding observation mask.

\paragraph{Conditional model.}
MotionSHAP-VAEAC uses three Transformer modules.  The full encoder
$q_\phi(\mathbf{z}\mid \mathbf{x},\mathbf{m})$ receives the complete sequence and mask and returns a diagonal Gaussian over per-frame latents,
\[
q_\phi(\mathbf{z}\mid \mathbf{x},\mathbf{m})
=
\mathcal{N}\!\left(\boldsymbol{\mu}_\phi,\operatorname{diag}(\boldsymbol{\sigma}^2_\phi)\right),
\quad
\mathbf{z}\in\mathbb{R}^{T\times d_z}.
\]
This encoder is used only during training.  The conditional prior
$p_\psi(\mathbf{z}\mid \mathbf{x}\odot\mathbf{m},\mathbf{m})$ sees only observed values and the mask and returns
$\mathcal{N}(\boldsymbol{\mu}_\psi,\operatorname{diag}(\boldsymbol{\sigma}^2_\psi))$.  The decoder $p_\theta(\hat{\mathbf{x}}\mid \mathbf{z},\mathbf{x}\odot\mathbf{m},\mathbf{m})$ maps a latent sample, observed values, and the mask to a complete sequence distribution.  Thus the prior and decoder form the conditional imputer used at explanation time.

\paragraph{Training.}
For each training sequence, masks are sampled from a mixture of temporal-window, spatial-joint, and element-wise masking distributions, ensuring support over the coalitions used at evaluation.  The VAEAC objective is
\begin{equation}
  \mathcal{L}_{\mathrm{VAE}}
  =
  \mathcal{L}_{\mathrm{rec}}
  +
  \beta(e)\,\mathcal{L}_{\mathrm{KL}}^{\mathrm{fb}}
  +
  \lambda_{\mathrm{reg}}\,\mathcal{L}_{\mathrm{reg}},
  \label{eq:loss_total}
\end{equation}
where $\beta(e)$ is annealed during training.  The reconstruction term is evaluated only on hidden, non-padded entries:
\begin{equation}
  \mathcal{L}_{\mathrm{rec}}
  =
  \mathbb{E}_{\mathbf{z}\sim q_\phi}
  \left[
    \frac{1}{|\bar{\Omega}|}
    \sum_{(j,f,t)\in\bar{\Omega}}
    \left(\hat{x}_{jft}-x_{jft}\right)^2
  \right],
  \label{eq:loss_rc}
\end{equation}
with $\bar{\Omega}=\{(j,f,t):m_{jft}=0\}$ intersected with valid frames.  The KL term is the closed-form diagonal Gaussian divergence
\begin{equation}
  \mathcal{L}_{\mathrm{KL}}
  =
  D_{\mathrm{KL}}\!\left(
  q_\phi(\mathbf{z}\mid\mathbf{x},\mathbf{m})
  \,\middle\|\,
  p_\psi(\mathbf{z}\mid\mathbf{x}\odot\mathbf{m},\mathbf{m})
  \right),
  \label{eq:loss_kl}
\end{equation}
optionally lower-bounded per latent dimension by free bits before averaging.  Following the stabilisation proposed for VAEAC~\citep{ivanov2018variational,olsen2022using}, $\mathcal{L}_{\mathrm{reg}}$ places weak normal and gamma priors on the prior encoder's mean and scale parameters.

\paragraph{Coalition imputation.}
For a coalition mask $\mathbf{m}_S$, MotionSHAP-VAEAC samples
$\mathbf{z}^{(r)}\sim p_\psi(\mathbf{z}\mid \mathbf{x}\odot\mathbf{m}_S,\mathbf{m}_S)$ and decodes
$\hat{\mathbf{x}}^{(r)}\sim p_\theta(\cdot\mid \mathbf{z}^{(r)},\mathbf{x}\odot\mathbf{m}_S,\mathbf{m}_S)$.  The hidden part of this completion supplies $\hat{\mathbf{x}}_{\bar S}^{(r)}$ in Eq.~\eqref{eq:vfn_mc}.

% ------------------------------------------------------------------
\subsection{MotionSHAP-Flow}
\label{sec:flow_imputer}
% ------------------------------------------------------------------

Our second imputer learns a conditional flow-matching velocity field~\citep{liu2022flow,lipman2022flow,tong2023conditional}.  It uses the same coalition mask $\mathbf{m}$ and observed values $\mathbf{x}\odot\mathbf{m}$, but represents the conditional distribution through an ODE rather than a latent-variable decoder.

\paragraph{Conditional model.}
Let $\mathbf{x}_1\sim p_{\mathrm{data}}$ and $\mathbf{x}_0\sim\mathcal{N}(\mathbf{0},\mathbf{I})$.  For $t\in[0,1]$, the conditional optimal-transport path is the affine bridge
\begin{equation}
  \mathbf{x}_t=(1-t)\mathbf{x}_0+t\mathbf{x}_1,
  \qquad
  \mathbf{u}_t=\frac{\mathrm{d}\mathbf{x}_t}{\mathrm{d}t}
  =\mathbf{x}_1-\mathbf{x}_0.
  \label{eq:condot_path}
\end{equation}
The velocity network $v_\eta(\mathbf{x}_t,t,\mathbf{m},\mathbf{x}_{\mathrm{obs}})$ is a Transformer over frame tokens.  Each token contains the current state $\mathbf{x}_t$, the observation mask $\mathbf{m}$, observed values $\mathbf{x}_{\mathrm{obs}}=\mathbf{x}_1\odot\mathbf{m}$, and a continuous embedding of flow time $t$.

\paragraph{Training.}
The flow objective regresses the network velocity to the analytic path velocity:
\begin{equation}
  \mathcal{L}_{\mathrm{flow}}
  =
  \mathbb{E}_{\mathbf{x}_1,\mathbf{x}_0,t,\mathbf{m}}
  \left[
    \frac{1}{|\bar{\Omega}|}
    \sum_{(j,f,t')\in\bar{\Omega}}
    \left\|
      v_\eta(\mathbf{x}_t,t,\mathbf{m},\mathbf{x}_{\mathrm{obs}})_{jft'}
      -
      (\mathbf{x}_1-\mathbf{x}_0)_{jft'}
    \right\|_2^2
  \right].
  \label{eq:flow_loss}
\end{equation}
Masks are sampled from the same temporal, spatial, and element-wise families used for VAEAC.  With small probability, all observations are dropped during training, giving the model an unconditional branch that can be used for classifier-free guidance.

\paragraph{Coalition imputation.}
Given a coalition mask $\mathbf{m}_S$, MotionSHAP-Flow samples completions by integrating the ODE
\begin{equation}
  \frac{\mathrm{d}\mathbf{x}_t}{\mathrm{d}t}
  =
  v_\eta(\mathbf{x}_t,t,\mathbf{m}_S,\mathbf{x}\odot\mathbf{m}_S),
  \qquad
  \mathbf{x}_0\sim\mathcal{N}(\mathbf{0},\mathbf{I}),
  \label{eq:flow_ode}
\end{equation}
from $t=0$ to $t=1$.  To keep known coordinates fixed while integrating hidden coordinates, we use RePaint-style harmonisation~\citep{lugmayr2022repaint}: after each numerical step, observed entries are reset to their location on the analytic bridge between the sampled noise and the target observation,
\begin{equation}
  \mathbf{x}_t \leftarrow
  \mathbf{m}_S\odot\big((1-t)\mathbf{x}_0+t\mathbf{x}\big)
  +(1-\mathbf{m}_S)\odot\mathbf{x}_t.
  \label{eq:repaint_harmonize}
\end{equation}
The resulting terminal state supplies $\hat{\mathbf{x}}_{\bar S}$ in Eq.~\eqref{eq:vfn_mc}.  Thus MotionSHAP-VAEAC and MotionSHAP-Flow instantiate the same cooperative game and Kernel SHAP estimator; they differ only in the learned conditional imputer.

% ------------------------------------------------------------------
\subsection{Baselines}
\label{sec:baselines}
% ------------------------------------------------------------------

All methods use identical coalitions and the same weighted least-squares Shapley estimator.  We compare the two on-manifold imputers above against three off-manifold baselines~\citep{frye2021shapley}: zero imputation, per-coordinate training-mean imputation, and marginal imputation by random training-sequence donors.  These baselines differ only in the distribution used to fill $\bar S$ in Eq.~\eqref{eq:motion_von}.

% ------------------------------------------------------------------
\subsection{Synthetic temporal benchmarks}
\label{sec:synthetic-benchmarks}
% ------------------------------------------------------------------

Real clinical evaluation is necessarily indirect: deletion, insertion, and related faithfulness scores are affected by classifier confidence and by how disruptive the perturbation is.  We therefore introduce two controlled motion benchmarks for which the data distribution and the oracle conditional $p(\mathbf{x}_{\bar S}\mid\mathbf{x}_S)$ are known.  This permits direct comparison of learned imputers against ground-truth Shapley values under temporal players with $K\in\{4,8\}$ and spatial players with $J=17$.

\paragraph{Gaussian motion.}
The Gaussian benchmark draws
\begin{equation}
  \mathbf{x}\sim
  \mathcal{N}\!\left(
    \mathbf{0},
    \Sigma_{\mathrm{joint}}\otimes I_F\otimes\Sigma_{\mathrm{time}}
  \right),
  \label{eq:gaussian_motion}
\end{equation}
equivalently
\[
\mathrm{Cov}(x_{jft},x_{j'f't'})
=
(\Sigma_{\mathrm{joint}})_{jj'}\,\delta_{ff'}\,(\Sigma_{\mathrm{time}})_{tt'}.
\]
By default, $\Sigma_{\mathrm{joint}}$ is equicorrelated with unit diagonal and off-diagonal $\rho$, and $\Sigma_{\mathrm{time}}$ is AR(1), $(\Sigma_{\mathrm{time}})_{tt'}=\alpha^{|t-t'|}$.  We also consider structured covariances estimated from real motion to encode skeleton adjacency and temporal persistence while preserving the Kronecker form.  For any coalition, the Gaussian conditional distribution is available in closed form, giving an oracle imputer for Eq.~\eqref{eq:motion_von}.

\paragraph{Burr copula motion.}
The Burr benchmark keeps the same Gaussian copula dependence but replaces Gaussian marginals with symmetric heavy-tailed Burr XII marginals.  Let $\mathbf{z}$ be drawn from Eq.~\eqref{eq:gaussian_motion}.  With Burr parameters $c,k>0$, standard normal cdf $\Phi$, and Burr quantile $Q_{\mathrm{Burr}}$, define
\begin{equation}
  x_{jft}
  =
  \operatorname{sign}(z_{jft})\,
  Q_{\mathrm{Burr}}\!\left(2\Phi(|z_{jft}|)-1;\,c,k\right).
  \label{eq:burr_motion}
\end{equation}
The default $c=k=2$ yields symmetric unit-scale heavy tails with tail index $ck=4$.  Oracle conditionals are obtained by inverting Eq.~\eqref{eq:burr_motion} to latent Gaussian space, applying the Gaussian conditional, and mapping samples back through the Burr transform.

\paragraph{Labels and black-box classifier.}
Labels follow the continuous interaction construction of Olsen et al.~\citep{olsen2022using}.  For temporal players, let $w_k(\mathbf{x})$ be the mean of $\mathbf{x}$ over joints, coordinates, and frames in window $k$, and set $u_k=\Phi(w_k/\sigma_k)$ using calibration standard deviations $\sigma_k$.  For each block of four windows, independently drawn coefficients $(a_r,b_r,c_r)$ define
\begin{equation}
  s_r(\mathbf{x})
  =
  a_r\sin(\pi u_{4r}u_{4r+1})
  +
  b_r u_{4r+2}\exp(c_r u_{4r+2}u_{4r+3}),
  \label{eq:olsen_score}
\end{equation}
and the score is $s(\mathbf{x})=\sum_r s_r(\mathbf{x})$.  The spatial benchmark applies the same four-variable interaction to four designated signal joints, replacing window means by joint means.  Class labels are the tertiles of $s(\mathbf{x})$ on a calibration sample, producing three approximately balanced classes.  The explained model is a small nonlinear classifier trained on the corresponding $K$ window means or $J$ joint means, so accurate attribution requires both modelling feature dependence and resolving the nonlinear interaction structure.